% ============================================================
%  《设计说明书》正文（精简版）：一 ~ 十二章
%  插图只保留 figures/ 中实际出图的 18 张；图号由 \figph 自动连续编号。
%  绘图语法与取图步骤见 chapters/figure_spec.tex（不编入 PDF）。
% ============================================================

\chapter{项目背景}

视觉语言导航要求智能体依据自然语言指令，在连续感知环境中自主行走到目标，是具身智能连接语言理解、环境感知与机器人决策的关键任务。连续环境视觉语言导航（VLN-CE）取消了离散导航图，智能体逐帧输出前进、转向、停止等底层动作；基于视频视觉语言模型的流式方法进一步以连续 RGB 观测为输入，通过滑动窗口与周期性慢更新记忆压缩历史信息，在有限上下文预算下处理长时间视频流。

围绕这类模型的工程化评测长期停留在脚本阶段，存在四个突出问题：一是评测入口分散，数据准备、切分生成、推理启动、指标统计各为独立脚本，实验间难以对齐；二是过程不可见，数小时的评测只能查看滚动日志，无法观察进度、显存、上下文长度与单步延迟；三是结果不可诊断，命令行只给出汇总指标，失败原因、轨迹形态与记忆状态需另行分析；四是复现成本高，配置、随机种子与中间结果未固化。多模态与时空记忆方向的实验还需要按“模态组合 × 记忆策略”成组对照，脚本方式下配置漂移与结果错配几乎不可避免。

本软件面向上述需求，把评测工作组织为“数据与配置—任务调度—推理回放—分析诊断”的完整闭环，并提供单局实时仿真，让模型轨迹、oracle 最短路径与人类示范轨迹同屏对照，把评测数字与现场行为对应起来。软件以浏览器为统一界面、以服务接口为统一数据通道，不训练模型、不分发数据集与权重，而是通过适配器协议对接外部推理后端，通过 Habitat 仿真器复用场景渲染与最短路径计算。

\chapter{软件功能介绍}

本软件按业务边界划分为八个功能区，见表~\ref{tab:functiongroups}。

\begin{navtable}{p{2.6cm}p{11.6cm}}{功能区与主要能力}{tab:functiongroups}
功能区 & 主要能力 \\
\midrule
总览 & 资产数量、任务状态、队列与指标趋势  \\
数据与配置 & 数据集登记/扫描、Episode 检索、观测时间轴、质量校验、三类切分、因果切片导出、配置模板 \\
模型档案 & 登记/编辑、可用性探测、自动配置  \\
评测执行 & 任务创建与配置快照、队列/并发/GPU 互斥、SSE 监控、暂停/恢复/取消/重跑、双适配器 \\
推理与回放 & 逐帧元数据采集、四区联动回放、记忆窗口与空间裁剪可视化 \\
分析诊断 & 指标对比、消融矩阵、配对 bootstrap、失败归因、长时曲线、报告导出 \\
交互式仿真 & 单局实时仿真、三种驾驶方式、俯视三线对照、画中画、录制导出 \\
资产与系统 & 备份库五类资产自动导入、用户/角色/审计、系统设置与 GPU 状态 \\
\end{navtable}

\section{总览}

首页汇总数据集、模型档案、评测任务与逐条结果的数量和状态，展示最近评测任务、本机存储与算力，并提供评测流程快捷入口。任务列表默认按创建时间最新在前。

\figph{ui01}{ui_01_overview.png}{系统首页总览}{界面截图：统计卡、队列、最近任务与指标趋势}

\section{数据与配置管理}

数据集以登记方式纳入管理，只记录名称、版本、类型与本地路径而不复制数据；扫描时解析 episode 清单，统计场景数、指令条数与轨迹长度并给出 ready / missing / invalid 状态。多源观测时间轴将 RGB、深度、位姿、动作按时间戳对齐展示并标出偏差区段；质量校验覆盖缺帧、时间戳乱序、位姿跳变、深度全零与空指令。切分支持全集、冒烟集、子集三种方式并保存生成条件与随机种子，因果切片导出只保留严格时间前缀的观测序列。配置模板预置“模态组合 × 记忆策略”组合供任务引用。

\figph{ui02}{ui_02_dataset_list.png}{数据集登记与扫描}{界面截图：数据集列表与扫描统计}
\figph{ui03}{ui_03_observation_timeline.png}{数据集详细信息查看}{界面截图：RGB / 深度 / 位姿 / 动作按时间戳对齐}

\section{评测任务调度}

任务创建向导依次选择数据集、切分、模型档案与配置模板，可覆盖窗口长度、慢更新间隔、裁剪比例、随机种子、并发数与 GPU 编号；提交时固化配置快照并执行因果性预检。任务按优先级排队，限制并行数并对同一 GPU 互斥；运行期以事件流推送进度、资源采样与日志，页面同步刷新进度条与曲线。任务控制支持取消、暂停、恢复与仅重跑失败项，暂停在 episode 边界生效，每条结果独立落盘以便续跑。

\figph{ui06}{ui_06_task_wizard.png}{新建评测任务向导}{界面截图：数据集/切分/模型/模板选择与参数覆盖}
\figph{ui07}{ui_07_task_monitor.png}{任务运行监控}{界面截图：进度、当前 episode、实时日志与资源曲线}

\section{推理与回放}

模型档案登记代码路径、权重路径与启动入口，支持可用性探测与从本地权重目录自动配置。软件以适配器协议对接外部推理后端，内置面向流式视觉语言导航模型的子进程适配器与离线演示适配器；调度真实模型时请求逐帧落盘，采集位姿、动作、单步时延、记忆窗口与裁剪前后 token 数，上游不支持时自动降级为仅汇总指标。回放工作台由俯视轨迹、第一视角画面、动作时间轴与记忆状态面板四区联动，支持拖动定位、单步、倍速与循环播放，并可叠加可通行区域掩膜与上下文长度曲线。

\figph{ui09}{ui_09_replay_workbench.png}{轨迹回放工作台}{界面截图：俯视轨迹、第一视角、运行状态与动作时间轴四区联动}

\section{指标分析与诊断}

软件聚合成功率、路径长度加权成功率、导航误差、oracle 成功率、归一化动态时间规整与轨迹长度，支持宏平均、微平均与按场景加权；提供多任务对比矩阵与基线差值列、按“模态组合 × 记忆策略”聚合的消融矩阵、同一 episode 集合上的配对 bootstrap 显著性检验、按步数分桶的长时性能曲线，以及超时、碰撞、偏离轨迹、打转、早停、不可达、解析错误、显存不足等九类失败归因。报告中心可导出含任务摘要、指标表与典型失败案例的 Markdown / CSV 报告。

\figph{ui12}{ui_12_analysis_compare.png}{多任务指标对比}{界面截图：多任务指标矩阵与相对基线的差值列}
\figph[1.0cm]{ui14}{ui_14_bootstrap.png}{配对 bootstrap 结果}{界面截图：均值差、置信区间与显著性标记}

\section{交互式仿真}

用户从已登记 episode 中选择一局，软件取其起点、目标与场景启动单局仿真。第一视角以 MJPEG 实时推送，俯视图同时绘制实际轨迹、人类示范轨迹与 oracle 最短路径；三种驾驶方式为手动按钮/键盘、模型驾驶与 oracle 自动走，可随时切换。读数包含步数、行走距离、测地距离、导航误差与 SPL 估计；第三视角画中画与机器狗模型仅用于可视化，不进入模型输入。会话可落盘帧图像并在停止后异步编码为 mp4。

\figph{ui18}{ui_18_sim_workbench.png}{交互式仿真工作台}{界面截图：第一视角 MJPEG、俯视轨迹、控制按钮与读数}
\figph[1.0cm]{ui19}{ui_19_sim_policy.png}{模型驾驶面板}{界面截图：模型加载、启用/暂停/停止与显存提示}

\section{资产同步}

绑定本地备份仓库后，软件按数据集、场景资源、模型档案、历史评测结果、其他资料五类扫描候选对象，提供只扫描不写库的预览模式；导入以来源键与内容签名幂等去重，不覆盖手工登记数据。自动同步由后台守护按间隔执行，开关与间隔保存在系统设置中，修改即时生效；仍在写入的结果文件保持“运行中”并按条数刷新进度。

\figph{ui23}{ui_23_backup_import.png}{备份库导入与自动同步}{界面截图：绑定目录、五类候选数、扫描预览与同步开关}

\section{系统管理}

提供用户与角色管理、操作审计、系统设置与运行状态监测。管理员、操作员、访客分别对应全部权限、读写任务与数据、只读查看；系统状态页汇总磁盘占用、数据库大小、队列长度与 GPU 显存，审计日志记录任务创建与取消、数据删除、配置修改、导入执行等关键操作。

\chapter{软件特性介绍}

\begin{description}[leftmargin=2.2em,itemsep=0.5em,font=\bfseries]
\item[一、评测流程闭环化。] 数据登记、切分生成、任务下发、过程监控、结果解析与报告导出在同一界面完成；配置在提交时固化为快照，历史任务不因模板变更而失去解释依据。

\item[二、过程可观测。] 运行期以事件流推送进度与资源采样，显存、延迟、上下文长度随时间成图，长时评测不再依赖滚动日志判断健康度。

\item[三、结果可诊断。] 保留每条 episode 的轨迹、动作序列、逐帧延迟与记忆窗口状态，可从“指标下降”下钻到“哪一条、哪一步、为什么失败”，人工修正后的归因可沉淀为案例。

\item[四、评测与现场一致。] 交互式仿真使用与评测相同的场景、动作空间与位姿口径，模型驾驶接入的正是被评测权重，oracle 最短路径与人类示范轨迹提供两条参照线。

\item[五、离线可演示。] 内置离线演示适配器读取预置结果，无 GPU、无数据集即可走通建库、任务、监控、分析、回放与报告全流程。

\item[六、模型无关与轻量部署。] 适配器协议隔离具体模型实现；SQLite 加单进程启动，前端构建产物由服务端托管，交付只需一条命令，无需数据库或消息中间件。

\item[七、因果性可验证、结果可复现。] 配置层输入白名单与运行层观测范围双重校验保证不泄漏未来信息；任务记录随机种子、配置快照、代码入口与环境摘要，切分保存生成条件，历史结果可复算。
\end{description}

\chapter{软件的运行环境介绍}

运行环境要求见表~\ref{tab:envreq}。离线演示与结果分析无需 GPU；接入真实推理后端或启用交互式仿真时建议配备 NVIDIA GPU 及对应驱动与 CUDA 环境，并与评测分卡使用以避免显存争用。交互式仿真所需的 \code{habitat_sim} 0.2.4、OpenCV 与模型驾驶所需的量化推理依赖均为可选，缺失时只影响对应页面。

\begin{navtable}{p{3.2cm}p{11.0cm}}{运行环境要求}{tab:envreq}
项 & 要求 \\
\midrule
CPU / 内存 & x86\_64 或 ARM64，Intel Core i5 及以上；离线模式不少于 8GB 内存 \\
磁盘 & 可用空间不少于 20GB（不含数据集、场景与模型权重） \\
GPU & 可选；真实推理与 Habitat 渲染建议单卡显存不少于 24GB \\
操作系统 & Ubuntu 20.04 及以上、Windows 10/11 \\
服务端 & Python 3.9+，FastAPI、Uvicorn、SQLAlchemy 2.0、Pydantic v2、PyYAML、SQLite 3（WAL） \\
前端 & Vue 3 / Vite 5 构建的 \code{frontend/dist} 由后端托管，运行时不需要 Node.js \\
浏览器 & Chrome / Edge / Firefox 近两年版本，1440×900 及以上分辨率 \\
\end{navtable}

部署采用单入口启动，首次启动自动建库、执行表迁移并写入种子数据：

\begin{lstlisting}[style=navplain]
python run.py --no-browser          # 默认监听 0.0.0.0:8080
\end{lstlisting}

应用首页为 \code{http://127.0.0.1:8080/}，接口文档为 \code{/docs}，健康检查为 \code{/health}。配置优先级为“环境变量 \textgreater{} \code{config/default.yaml} \textgreater{} 内置默认值”；其中备份库自动同步开关等运行期参数保存在 \code{setting} 表，修改后立即生效。

\chapter{系统总结构}

\section{总体架构}

本软件采用前后端分离、单机一体化部署结构，由浏览器前端、接口服务层、业务服务层、仿真子系统、适配器层、数据存储层与外部推理/仿真后端七部分组成，如图~\ref{fig:arch01} 所示；各层职责见表~\ref{tab:layers}。

\figph[1.0cm]{arch01}{arch_01_system_overview.png}{系统总体架构图}{PlantUML 组件图：七层结构与层间调用关系}

\begin{navtable}{p{3.0cm}p{8.0cm}p{3.2cm}}{分层职责}{tab:layers}
层 & 职责 & 主要技术 \\
\midrule
浏览器前端层 & 页面渲染与交互；HTTP 取数、事件流接收状态、Canvas 绘轨迹、MJPEG 消费帧流 & Vue 3 / Pinia / ECharts / Canvas \\
接口服务层 & 路由分发、参数校验、身份认证、统一响应封装、路径白名单与令牌查询参数通道 & FastAPI / Pydantic \\
业务服务层 & 数据扫描、切分、质量校验、任务编排、指标聚合、失败归因、回放解析、报告与导入 & Python 服务模块 \\
仿真层 & 场景构建、引擎循环、会话管理、自动策略、录制与俯视绘制；引擎线程独占渲染上下文 & habitat\_sim / OpenCV \\
适配层 & 统一准备/执行/汇总接口；生成后端配置、启动子进程、增量解析输出与逐帧元数据 & Adapter 协议 \\
数据访问与存储层 & SQLite 结构化数据 + 工作目录文件（快照、日志、逐帧文件、帧图像、录制、报告） & SQLite(WAL) / 本地文件 \\
外部支撑层 & 视觉语言推理、动作生成与场景渲染；本软件不实现训练与推理计算 & StreamVLN / Habitat \\
\end{navtable}

数据存储层由 SQLite 数据库与工作目录组成：数据库保存数据集、切分、模型档案、任务、逐条结果、资源采样、产物索引、导入账本、用户与审计等信息；工作目录保存配置快照、运行日志、逐帧元数据、帧图像、仿真录制与报告文件，大文件不入库，仅登记路径。

\section{评测任务执行总流程}

一次完整评测的执行过程如图~\ref{fig:flow01} 所示：用户选择数据集、切分、模型与配置模板并提交，服务端校验参数与因果性后写入任务记录并进入队列；调度器在并发额度与 GPU 互斥条件满足时取出任务，生成配置快照并启动适配器；适配器拉起子进程逐条执行 episode，将进度、资源采样、逐帧元数据与日志以标准格式输出；服务端按行解析并写库，同时向前端推送事件；任务结束后聚合指标、生成回放数据与报告索引。

\figph[1.0cm]{flow01}{flow_01_eval_task.png}{评测任务执行总流程}{PlantUML 活动图：从任务提交到结果聚合的跨层执行流程}

\section{技术栈与设计约束}

后端采用 Python 3.9+ / FastAPI / SQLAlchemy 2.0 / Pydantic v2，前端采用 Vue 3 / Vite 5 / Element Plus / ECharts / Canvas，调度为自研轻量队列（asyncio + 子进程），不依赖 Redis 或 Celery。设计约束包括：离线可演示、不含数据与权重、因果性可验证、单机可交付、模型无关、仿真可选。

\chapter{系统各子系统介绍}

\section{数据与配置管理子系统}

该子系统解决“评测输入从哪来、是否可用、如何固定”的问题。数据集登记后解析 episode 清单写入索引表，形成可按场景、指令长度、轨迹长度检索的数据视图；观测时间轴按时间戳对齐多源数据并标记偏差区段；质量校验输出问题清单；切分按全集、冒烟集、子集三种方式生成并保存条件与种子；因果切片导出截取至指定步后二次校验不含未来帧。处理流程见图~\ref{fig:flow02}。

\figph[1.0cm]{flow02}{flow_02_data_manage.png}{数据与配置管理子系统流程}{PlantUML 活动图：登记 → 扫描 → 校验 → 切分 → 因果切片导出}

\section{评测任务调度子系统}

任务提交时把模板参数与覆盖项合并为配置快照，此后模板变更不影响该任务。调度器按优先级出队，以并发额度控制同时运行的任务数，同一编号 GPU 上的任务互斥；任务启动时创建工作目录、写入配置快照与启动命令并记录进程号。运行期以事件流推送进度、资源采样与日志，进度事件节流，日志同时落盘。取消时终止子进程，暂停在 episode 边界生效，重试仅针对失败 episode 且保留已完成结果。流程见图~\ref{fig:flow03}。

\figph[1.0cm]{flow03}{flow_03_scheduler.png}{任务调度子系统流程}{PlantUML 活动图：入队 → 并发/GPU 判定 → 启动 → 事件回写 → 结束与重试}

\section{推理与回放子系统}

适配器以统一接口封装差异：准备阶段校验路径、生成后端配置、创建输出目录；执行阶段逐条调用后端并收集标准记录；汇总阶段统计指标并落盘。子进程适配器按约定命令启动后端并增量读取结果，任务开始时向上游请求逐帧元数据；离线演示适配器读取预置结果并按种子加入可复现扰动。

逐帧元数据是回放的数据基础，每步记录位姿、朝向、动作、单步时延、窗口帧数、是否慢更新以及裁剪前后 token 数，以行为单位追加落盘，适配器按帧序规整后写入任务目录；上游不支持该能力时记录日志并降级为仅汇总指标。回放工作台四区同步刷新，可直观看到某一步动作之后记忆窗口的变化。完整流程见图~\ref{fig:flow04}。

\figph[1.0cm]{flow04}{flow_04_adapter_replay.png}{推理适配器与逐帧回放流程}{PlantUML 活动图：prepare → 子进程逐帧落盘 → 解析规整 → 回放接口 → 四区工作台}

\section{指标分析与诊断子系统}

指标聚合按任务读取逐条结果，计算 SR、SPL、NE、OS、nDTW 与平均轨迹长度并给出宏/微平均；多任务对比生成指标矩阵与基线差值列，消融矩阵按“模态组合 × 记忆策略”两维聚合，配对 bootstrap 在同一 episode 集合上重采样给出均值差与 95\% 置信区间。失败归因按规则自动归类并允许人工复核，长时性能曲线按步数分桶观察退化趋势。处理流程见图~\ref{fig:flow05}。

\figph[1.0cm]{flow05}{flow_05_analysis.png}{指标分析与诊断流程}{PlantUML 活动图：逐条结果 → 指标聚合 → 对比/消融/显著性 → 失败归因 → 报告导出}

\section{交互式仿真子系统}

用户选定 episode 后读取场景、起点、目标与参考路径，在 Habitat 中构建环境并启动引擎线程；引擎按固定节拍应用动作、渲染第一视角、绘制俯视底图与三条轨迹线，整帧只编码一次 JPEG 供多观众共享。手动模式接受按钮与键盘指令，oracle 模式沿 navmesh 最短路径逐点跟随，模型模式把当前观测送入推理线程并通过动作队列与引擎解耦，三者可随时切换。诊断读数与评测口径一致，画中画与机器狗仅用于可视化。流程见图~\ref{fig:flow06}。

\figph[1.0cm]{flow06}{flow_06_sim.png}{交互式仿真子系统流程}{PlantUML 活动图：会话启动 → 引擎循环 → 三种驾驶方式 → 帧分发 → 录制导出 → 释放}

\section{资产同步子系统}

软件按五类资产扫描备份仓库：数据集解析为 episode 索引，场景与资料登记为文件索引，模型目录补齐模型档案的路径与启动入口，历史结果还原为可回放任务。导入以来源键与内容签名幂等判定，预览模式只扫描不写库；自动同步开关与间隔保存在系统设置中。仍在写入的结果文件保持“运行中”并按条数刷新进度，支持“边跑边看”。流程见图~\ref{fig:flow07}。

\figph[1.0cm]{flow07}{flow_07_backup_import.png}{资产同步导入流程}{PlantUML 活动图：扫描 → 分类 → 签名比对 → 预览/写入 → 增量同步}

\chapter{关键机制设计}

\section{适配器协议与逐帧元数据契约}

适配器协议把“平台如何组织评测”与“模型如何推理”解耦，统一约定 \code{prepare}、\code{run}、\code{finish} 三个阶段。子进程适配器拼接启动命令并追加逐帧落盘参数，上游脚本不支持时在准备阶段探测能力并自动降级；离线演示适配器读取预置结果。逐帧文件每行一帧，字段固定，Mock 与真实适配器共用，见表~\ref{tab:framecontract}；解析对缺字段、半行与重复帧容错。

\begin{navtable}{p{4.2cm}p{2.2cm}p{7.8cm}}{逐帧元数据字段契约}{tab:framecontract}
字段 & 类型 & 说明 \\
\midrule
\code{idx} / \code{t} & int / float & 步序号（从 0 开始）、相对 episode 开始的秒数 \\
\code{pose} / \code{yaw} & 数组 / float & 世界坐标位置（米）与偏航角（弧度） \\
\code{action} & string & forward / left / right / stop \\
\code{latency_ms} & float & 该步决策耗时（毫秒） \\
\code{window_len} & int & 当前记忆/视觉窗口内帧数 \\
\code{slow_update} & bool & 该步是否发生慢更新（历史重采样） \\
\code{tokens_before_prune} / \code{tokens_after_prune} & int & 裁剪前后 memory token 数 \\
\code{ctx_len} & int & 当前 prompt token 数 \\
\end{navtable}

\section{配置快照与因果性双层校验}

任务提交时把模板参数、覆盖项、随机种子、数据集与切分标识、模型档案标识、代码入口与环境摘要合并为不可变配置快照写入任务记录，此后模板编辑不影响历史任务；切分独立保存生成条件与种子，可复算。因果性第一层在配置层校验输入白名单（如 \code{rgb}、\code{depth}、\code{pose}、\code{action}），非法字段直接拒绝并返回错误码 1006；第二层在运行层按时间前缀约束观测读取，因果切片导出前二次校验不含未来帧，越界断言触发即中断任务并写审计。

\section{任务队列、资源互斥与实时监控}

调度器以 asyncio 维护待执行队列，并发上限默认 2 且可在系统设置中即时调整，同一 GPU 编号上的任务互斥；任务在 episode 边界检查暂停与取消信号，保证不产生半条轨迹。资源采样周期性记录显存峰值、单步延迟、上下文长度与 token 数并写入数据库。服务端以 \code{text/event-stream} 推送 \code{progress}、\code{sample}、\code{log}、\code{state}、\code{done} 五类事件，进度节流不短于 500 毫秒，事件队列在消费端过慢时丢弃最旧事件；前端按事件类型分别更新进度条、资源曲线、日志面板与状态标签，代理不支持 SSE 时可轮询降级。

\section{仿真三线程解耦与单帧单编码}

引擎循环、会话管理与模型推理分置不同线程：引擎独占 headless GL 上下文并按节拍产出整帧，会话层维护动作队列与观看者列表，推理线程消费同一份第一视角观测并把动作写入队列。整帧只编码一次 JPEG，字节由多观众共享；画中画只切显隐、不进入模型输入。同一时刻只允许一个仿真会话，启动新会话先停旧的，空闲超时自动停止；停止时串行释放引擎线程、GL 上下文与显存。仿真调试禁用 \code{--reload}，避免进程重建强杀引擎线程导致 GL 上下文残留。

\section{幂等导入与显著性检验}

导入以 \code{source_key + signature} 幂等判定，重复执行只处理新增或变化对象；\code{dry_run} 只扫描不写库。外部结果文件在 \code{running_grace_s}（默认 600 秒）内仍有写入时保持“运行中”并按实际条数刷新进度，停止写入后自动转为已完成。配对 bootstrap 在同一 episode 集合上按观测差值中心化后做双侧检验，重采样默认 1000 次，输出均值差、95\% 置信区间与显著性标记，最小 p 值受重采样次数约束。

\section{鉴权双通道与可选依赖降级}

常规接口使用 \code{Authorization: Bearer <token>}；\code{<img>}、MJPEG 与文件下载无法携带自定义请求头，改用 \code{?token=<access_token>} 通道单独校验。所有 \code{path} 类参数必须落在工作目录白名单根内防止目录穿越，管理类接口按角色二次判定；用户不存在时仍执行一次哈希计算，避免通过响应时间枚举账号。\code{habitat_sim} 与推理框架、上游逐帧能力、原始 RGB 帧、掩膜文件、SSE 通道均为可选或可降级，缺失时只影响对应功能并给出明确提示。

\chapter{接口设计}

\section{通用约定与错误码}

接口前缀为 \code{/api}，编码 UTF-8，请求与响应为 JSON；统一响应封装为 \texttt{\{"code":0,"message":"ok","data":\{...\}\}}。分页参数为 \code{?page=1} 与 \code{page_size=20}，响应 \code{data} 含 \code{items}、\code{total}、\code{page}、\code{page_size}；所有 \code{path} 类参数执行白名单校验。错误码见表~\ref{tab:errorcode}。

\begin{navtable}{p{2.0cm}p{1.8cm}p{10.4cm}}{统一错误码}{tab:errorcode}
code & HTTP & 含义 \\
\midrule
0 & 200 & 成功 \\
1001 & 400 & 参数校验失败 \\
1002 & 401 & 未登录或 Token 失效 \\
1003 & 403 & 无权限 \\
1004 & 404 & 资源不存在 \\
1005 & 409 & 状态冲突（任务已在运行、仿真已有会话） \\
1006 & 422 & 因果性校验失败 \\
2001 & 500 & 内部错误 \\
2002 & 503 & 推理后端不可用（GPU/权重缺失） \\
\end{navtable}

\section{接口分组}

接口按业务域分为九组，如表~\ref{tab:apigroups} 所示。

\begin{navtable}{p{3.0cm}p{11.2cm}}{接口业务域与主要接口}{tab:apigroups}
业务域 & 主要接口 \\
\midrule
认证与用户 & \code{POST /api/auth/login}、\code{GET /api/users} \\
数据集 & \code{GET /api/datasets}、\code{POST /api/datasets/:id/scan} \\
质量与切分 & \code{POST /api/datasets/:id/validate}、\code{/splits}、\code{/export-causal} \\
模板与模型 & \code{GET /api/templates}、\code{/models}、\code{POST /api/models/:id/probe} \\
任务 & \code{POST /api/tasks}、\code{GET /api/tasks/:id/events}、\code{/pause}、\code{/retry-failed} \\
分析与回放 & \code{GET /api/analysis/compare}、\code{POST /api/analysis/bootstrap}、\code{GET /api/replay/:task/:row} \\
仿真 & \code{POST /api/sim/sessions}、\code{GET /api/sim/state}、\code{/stream}、\code{POST /api/sim/actions}、\code{/policy} \\
导入与报告 & \code{GET/POST /api/import/backup}、\code{POST /api/reports} \\
系统 & \code{GET /api/overview}、\code{/system/status}、\code{/settings}、\code{/audit} \\
\end{navtable}

\section{关键报文示例}

创建任务请求、响应与因果性预检失败示例：

\begin{lstlisting}[style=navplain]
POST /api/tasks
{"name":"vision_geo_state_seed42","dataset_id":1,"split_id":4,
 "model_profile_id":1,"template_id":4,
 "overrides":{"history_len":8,"prune_ratio":0.35,"seed":42}}
-> {"code":0,"data":{"id":12,"status":"queued",
   "progress_total":1839,"config_snapshot":{"...":"..."}}}
-> {"code":1006,
   "message":"输入 rgb_future 不在白名单 [rgb,depth,pose,action]"}
\end{lstlisting}

回放接口返回任务、结果行、指令分段、配置快照、逐帧数组、汇总、几何与记忆序列；无逐帧文件时帧数组为空、只显示汇总指标与参考路径并给出提示，无原始 RGB 时返回由帧元数据生成的示意帧 SVG。仿真状态快照包含会话状态、步数与距离、测地距离/NE/SPL 估计、位姿、三条轨迹线、控制开关、帧率与录制进度。

\section{权限与兼容性}

仿真启动、操控、停止、模型驾驶与录制导出限管理员与操作员；状态、帧流与文件下载走 \code{?token=} 通道；用户管理与审计限管理员。接口按相对路径版本演进，新增字段向后兼容，字段缺失时前端按缺省值降级显示。

\chapter{人机接口设计}

\section{布局与导航}

前端采用统一设计系统：固定左侧导航栏、顶部标题与副标题区、右侧内容区；颜色、字号、间距与图标由统一变量管理，浏览器标签、侧边栏与登录页共用同一应用图标。列表、表单、弹窗与状态标签复用 Element Plus 组件，图表使用 ECharts，轨迹与掩膜使用 Canvas 自绘。主要前端路由见表~\ref{tab:routes}；路由守卫统一处理未登录跳转、管理员页面角色判定与浏览器标题设置。

\begin{navtable}{p{4.4cm}p{4.2cm}p{5.6cm}}{主要前端路由}{tab:routes}
路由 & 页面 & 主要交互 \\
\midrule
\code{/login} & 登录 & 口令登录、跳转回原地址 \\
\code{/overview} & 总览 & 统计卡、最近任务、快捷入口 \\
\code{/datasets}、\code{/datasets/:id} & 数据集与详情 & 登记、扫描、检索、时间轴、校验、切分 \\
\code{/models} & 模型档案 & 登记、自动配置、可用性探测 \\
\code{/tasks}、\code{/tasks/:id} & 任务列表与监控 & 创建向导、队列、SSE 监控、任务控制 \\
\code{/analysis}、\code{/reports} & 分析与报告 & 对比、消融、显著性、归因、报告导出 \\
\code{/replay} & 轨迹回放 & 四区联动、逐帧、记忆与几何叠加 \\
\code{/sim} & 交互式仿真 & MJPEG、俯视三线、三种驾驶、录制 \\
\code{/backup-import} & 备份库导入 & 扫描预览、立即导入、自动同步 \\
\code{/settings}、\code{/users}、\code{/audit} & 系统管理 & 设置、用户、审计与状态 \\
\end{navtable}

\section{操作旅程与状态反馈}

主操作旅程为：登录 → 总览 → 登记/确认数据集与模型 → 生成切分 → 向导创建任务 → 监控页观察进度与资源曲线 → 分析页对比与归因 → 回放页定位失败步 → 报告中心导出材料；离线演示模式可跳过真实推理。表单校验失败在字段下方就地提示，异步任务用进度条与状态标签反馈，危险操作需二次确认。任务列表默认按创建时间最新在前；外部评测结果文件在写入期间显示“运行中”并按实际条数显示进度，停止写入后自动转为已完成；模型档案页提供“自动配置”，系统设置中的自动同步开关无需重启即时生效。仿真页键盘为 ↑/W 前进、←/A 左转、→/D 右转、空格停止，每次前进 0.25 米、转向 15 度。

\chapter{测试计划}

\section{测试范围与环境}

测试覆盖功能正确性、接口契约、因果性约束、容错降级、性能并发、安全与兼容性，不覆盖上游模型本身的精度与训练。离线环境无 GPU、无数据集，使用 Mock 适配器与内置示例结果验证全功能链路；在线环境为 Ubuntu + Python 3.9+ + habitat\_sim + NVIDIA GPU（本机 8×L40S），验证真实推理、逐帧回放与仿真。测试数据为 R2R-CE val\_unseen、冒烟切分 20 条与示例结果目录。当前版本以人工端到端验证为主，自动化测试框架为后续补充项。

\section{功能验收用例}

\begin{navtable}{p{1.6cm}p{13.6cm}}{功能验收用例}{tab:testcases}
编号 & 用例与判据 \\
\midrule
TC-01 & 首次启动自动建库，30 秒内可登录 \\
TC-02 & 登记示例数据集并扫描，episode 可检索、可看观测时间轴 \\
TC-03 & 生成 smoke 切分、导出因果切片，文件中无晚于截止步的帧 \\
TC-04 & Mock 任务 20 条约 5 分钟内跑完，监控页进度/日志/资源曲线正常 \\
TC-05 & 指标非空，两任务差值列正确；失败 episode 自动归因 \\
TC-06 & 回放页四区联动、拖动时间轴同步刷新；导出 Markdown 报告 \\
TC-07 & 真实任务落逐帧文件且行数等于步数；上游不支持时降级不报错 \\
TC-08 & 仿真三种驾驶可用、帧率不低于 15 fps；录制 mp4 帧数一致并释放显存 \\
\end{navtable}

\section{技术验收与当前状态}

技术验收包括：违规配置触发因果性校验并返回 1006；任务取消后子进程退出无残留；服务重启后遗留 running 任务被正确标记而外部导入任务不受影响；逐帧解析对缺字段、半行、重复帧容错；越权访问返回 401/403、路径穿越被白名单拦截、无令牌帧请求返回 401。性能判据为 1000 条 episode 聚合小于 1 秒、回放帧切换小于 100 毫秒、仿真不低于 15 fps、最多 2 个评测任务并行。已完成端到端验证覆盖建库/鉴权、数据与切分、任务调度与双适配器、分析与报告、轨迹回放、交互式仿真、备份库导入与真实适配器逐帧元数据采集；仓库暂无自动化测试，以人工回归记录结果。

\chapter{系统出错处理设计}

\section{错误分类}

软件把错误分为参数与校验错误、认证与授权错误、资源与状态冲突、依赖不可用、内部异常五类，分别映射到表~\ref{tab:errorcode} 的错误码。前端按 \code{code} 字段统一提示，业务层抛出带上下文的异常并记录日志，接口层转换为统一响应，不向前端暴露堆栈。

\section{分层容错与降级策略}

\begin{navtable}{p{4.6cm}p{9.6cm}}{降级策略}{tab:degrade}
场景 & 处理方式 \\
\midrule
上游不支持逐帧落盘 & 记录日志，降级为仅汇总指标，评测继续 \\
无逐帧文件 / 无原始 RGB & 回放显示汇总指标、参考路径或示意帧 SVG 并提示 \\
无可通行区域掩膜 & 掩膜接口返回 1004，前端隐藏叠加层 \\
SSE 在代理后失效 & 前端切换轮询降级 \\
\code{habitat_sim} 缺失 & 仅仿真页报“仿真组件不可用”，其余功能不受影响 \\
显存不足 / 权重或 GPU 缺失 & 模型驾驶回退手动或返回 2002 并给出缺失项 \\
SQLite 写锁 & WAL + 单写线程 + 批量提交，运行时不直接删库文件 \\
子进程崩溃 & 标记当前 episode 失败而非中断全任务，日志落盘 \\
\end{navtable}

\section{任务失败、日志与资源释放}

每条 episode 结果独立落盘，失败只标记该条并计入归因，任务继续；用户可仅重跑失败项，已完成结果保留。服务重启后由平台调度的 running 任务标记为失败并支持重跑失败项，外部导入结果不受影响。运行日志写入 \code{storage/logs/app.log}，任务日志写入 \code{storage/runs/<task>/run.log}；健康检查 \code{GET /health} 返回状态与版本，关键写操作同时写审计日志。故障定位按“任务编号 → 运行日志 → episode 结果 → 逐帧文件 → 回放画面”下钻。仿真会话空闲超过 \code{idle_timeout_s}（默认 900 秒）自动停止，停止时串行释放引擎线程、GL 上下文与显存；服务退出前按序停止后台守护与仿真会话。数据库启用 WAL 与外键约束，导入幂等，删除记录前检查任务引用，恢复时从 SQLite 备份与工作目录整体还原。

\chapter{系统维护设计}

\section{配置与数据库}

配置优先级为“环境变量 \textgreater{} \code{config/default.yaml} \textgreater{} 内置默认值”；修改 YAML 后需重启，而备份库自动同步开关、间隔与绑定目录等运行期参数保存在 \code{setting} 表，修改后立即生效。仿真参数（场景根目录、GPU、帧率、JPEG 质量、观看上限、空闲超时、显存阈值、画中画、录制、机器狗）集中在 \code{sim} 配置段。首次启动自动建库、执行表迁移并写入种子数据；版本由 \code{schema_version} 记录，迁移脚本按序执行，备份使用 SQLite 在线备份接口或停机后复制 \code{eval.db}。数据集、场景、模型与结果等大体积内容只登记路径，不随数据库搬迁。

\section{运行目录、监控与升级}

运行目录按职责划分：\code{storage/eval.db} 保存结构化数据，\code{storage/runs/} 保存配置快照、运行日志与逐帧文件，\code{storage/sim/} 保存录制帧与 mp4，\code{storage/reports/} 保存报告，\code{storage/logs/} 保存应用日志，\code{storage/cache/} 与 \code{storage/exports/} 保存缓存与导出文件；结果保留天数与录制轮转数量由系统设置控制，清理时按任务编号整体删除。运行监控依托健康检查、系统状态页与 GPU 状态接口。升级流程为：停止服务并确认无运行中任务 → 备份数据库与工作目录 → 替换后端代码与前端 dist → 启动自动迁移 → 健康检查与冒烟回归；异常时回滚代码并还原备份。仅更新前端 dist 时无需重启后端，修改后端代码或 YAML 才需要重启。

\section{已知限制与交付材料}

预留扩展点包括新增适配器实现、指标与归因规则、配置模板、资料索引类别与多机调度接口。当前限制：仿真需装有 \code{habitat_sim} 的环境与场景文件；数据库切分只用于结果对齐；逐帧回放完整度取决于上游是否支持逐帧落盘；平台调度任务断点续跑以 episode 为粒度；大规模训练数据默认只登记路径；仓库暂无自动化测试。

软件交付物包括可运行的源码目录与依赖清单、前端构建产物、配置样例、项目设计文档（需求、数据库、接口、专题设计）。维护时以 \code{docs/} 下文档为基线同步更新本说明书，插图按 \code{figures/} 中的命名维护。
